\section{Билет 38. Основные положения молекулярно-кинетической теории. Число Авогадро. Давление идеального газа на стенку с выводом.}

\begin{definition}
\textbf{Молекулярно-кинетическая теория (МКТ)} --- раздел физики, изучающий строение и свойства веществ на основе представлений о том, что все тела состоят из мельчайших частиц (молекул и атомов), находящихся в непрерывном хаотическом движении.
\par\noindent\textbf{Три основных положения МКТ:}
\begin{enumerate}
    \item Все вещества --- жидкие, твёрдые и газообразные --- образованы из мельчайших частиц: молекул, которые сами состоят из атомов. Молекулы и атомы электрически нейтральны.
    \item Атомы и молекулы находятся в непрерывном хаотическом тепловом движении.
    \item Частицы взаимодействуют друг с другом силами электрической природы; гравитационное взаимодействие между ними пренебрежимо мало.
\end{enumerate}
\end{definition}

\begin{definition}
\textbf{Число Авогадро} $N_A$ --- физическая постоянная, равная числу структурных элементов (атомов, молекул, ионов), содержащихся в одном моле любого вещества:
\begin{equation}
    N_A = 6.02 \cdot 10^{23} \text{ моль}^{-1}. \label{eq:avogadro_38}
\end{equation}
Количество вещества $\nu$ определяется как отношение числа частиц $N$ к числу Авогадро:
\begin{equation}
    \nu = \frac{N}{N_A}.
\end{equation}
Молярная масса $M$ связана с массой одной молекулы $m_0$ соотношением: $M = N_A m_0$.
\end{definition}

\begin{theorem}[Вывод давления идеального газа на стенку]
Рассмотрим идеальный газ в сосуде кубической формы с ребром $l$. Пусть в единице объёма содержится $n$ молекул (концентрация), каждая массой $m_0$.

\textbf{1. Импульс, передаваемый одной молекулой.}
Молекула, движущаяся перпендикулярно стенке со скоростью $v_x$, при абсолютно упругом ударе изменяет свою проекцию импульса на $2m_0 v_x$. По закону сохранения импульса, стенка получает импульс:
\begin{equation}
    \Delta p_x = 2 m_0 v_x.
\end{equation}

\textbf{2. Число ударов о стенку за время $\Delta t$.}
За время $\Delta t$ о стенку площадью $S$ успеют удариться только те молекулы, которые находятся в объёме $S \cdot v_x \Delta t$ и движутся к стенке. Учитывая, что движение хаотично, к данной стенке летит $1/6$ всех молекул (по трём осям, в двух направлениях):
\begin{equation}
    N_{\text{уд}} = \frac{1}{6} n S v_x \Delta t.
\end{equation}

\textbf{3. Сила давления.}
Общий импульс, переданный стенке за время $\Delta t$:
\begin{equation}
    \Delta P = N_{\text{уд}} \cdot 2 m_0 v_x = \frac{1}{6} n S v_x \Delta t \cdot 2 m_0 v_x = \frac{1}{3} n m_0 v_x^2 S \Delta t.
\end{equation}
Сила --- это изменение импульса в единицу времени:
\begin{equation}
    F = \frac{\Delta P}{\Delta t} = \frac{1}{3} n m_0 v_x^2 S.
\end{equation}

\textbf{4. Давление.}
Давление --- сила, действующая на единицу площади:
\begin{equation}
    p = \frac{F}{S} = \frac{1}{3} n m_0 \langle v_x^2 \rangle.
\end{equation}
Учитывая, что для хаотического движения $\langle v^2 \rangle = \langle v_x^2 \rangle + \langle v_y^2 \rangle + \langle v_z^2 \rangle = 3\langle v_x^2 \rangle$, получаем \textbf{основное уравнение МКТ}:
\begin{equation}
    p = \frac{1}{3} n m_0 \langle v^2 \rangle = \frac{2}{3} n \left\langle \frac{m_0 v^2}{2} \right\rangle = \frac{2}{3} n \langle E_{\text{пост}} \rangle. \label{eq:mkt_main_38}
\end{equation}
\end{theorem}

\begin{property}
Из основного уравнения МКТ \eqref{eq:mkt_main_38} следует, что:
\begin{enumerate}
    \item Давление идеального газа пропорционально концентрации молекул $n$ и средней кинетической энергии поступательного движения молекул $\langle E_{\text{пост}} \rangle$.
    \item Связь с уравнением состояния: подставляя $\langle E_{\text{пост}} \rangle = \frac{3}{2} kT$ (где $k$ --- постоянная Больцмана), получаем $p = n k T$, что эквивалентно уравнению Менделеева--Клапейрона $pV = \nu RT$ при учёте $n = N/V$ и $R = k N_A$.
\end{enumerate}
\end{property}

\begin{remark}
\textbf{Физический смысл температуры.} Сравнивая $p = \frac{2}{3} n \langle E_{\text{пост}} \rangle$ и $p = n k T$, получаем фундаментальное соотношение:
\begin{equation}
    \langle E_{\text{пост}} \rangle = \frac{3}{2} k T.
\end{equation}
Это означает, что абсолютная температура является мерой средней кинетической энергии поступательного движения молекул идеального газа. Температура --- статистическая величина, характеризующая не отдельную молекулу, а ансамбль частиц.
\end{remark}
